% ============================================================
% Frohmanian Tether Symbol - Recommended LaTeX Implementation
% Per user's detailed guidance (2026-06-03)
% ============================================================
%
% RECOMMENDED METHOD (Best practical choice for papers):
% Clean, maintainable, works with pdflatex and lualatex.
% Uses \mathfrak for a distinctive mathematical look.
%
% This is the sweet spot: easy + professional.
% Adjust \mkern for perfect visual spacing.

\newcommand{\FT}{\mathfrak{T}\mkern-1.5mu\mathfrak{F}}   % or use your actual drawn symbol
\newcommand{\FTh}{\FT}

% Usage examples:
% The core object is the \FT.
% The \FTh{} Theorem states that every smooth divergence-free initial
% velocity field on \mathbb{T}^3 generates a unique globally smooth solution.

% ------------------------------------------------------------
% ALTERNATIVE: Unicode approach (modern, if using unicode-math + lualatex)
% \usepackage{unicode-math}
% \newcommand{\FT}{𝔉𝕋}   % or find a better precomposed ligature

% ------------------------------------------------------------
% ALTERNATIVE: TikZ for exact hand-drawn match (use sparingly, for figures)
% \usepackage{tikz}
% \newcommand{\FT}{%
%   \begin{tikzpicture}[baseline=(char.base), scale=0.9]
%     \node[inner sep=0pt] (char) at (0,0) {$\mathfrak{T}$};
%     \node[inner sep=0pt, xshift=0.8pt] at (char.center) {$\mathfrak{F}$};
%   \end{tikzpicture}%
% }

% ------------------------------------------------------------
% Why this recommendation:
% - Clean and easy to maintain
% - Works in standard LaTeX setups
% - You control the spacing
% - Easy to evolve if a true custom font glyph becomes available later
% - Matches the "Best Practical Method" in the provided LaTeX ligature techniques.

% For the absolute best visual (professional font ligature), use fontspec + lualatex
% with a custom stylistic set or font feature. Only pursue if typography is critical
% and you have font design resources.

% This preamble snippet should be included in the main LaTeX manuscript.
% See also: Frohmanian_Tether_Naming_Symbol_Standard.md (copied locally).
